\section{The \texttt{catena.numbers} module}
\label{sec:numbers}

The \texttt{catena.numbers} module provides two sub-modules:
\texttt{constants}, which pre-constructs SCF representations of well-known
mathematical constants, and \texttt{randoms}, which supplies generators
and SCF factories whose coefficients are drawn from the Gauss--Kuzmin
distribution.  The present section describes both.  The mathematical
justification for the random generators requires results on the measure
theory of continued fractions that go beyond the scope of this version of
the document; the relevant facts are stated without proof and deferred to
a forthcoming section.

\subsection{Constants}
\label{subsec:numbers-constants}

The module exposes the following pre-built SCF objects.

\begin{description}[leftmargin=2em, style=nextline]
  \item[$\mathtt{e}$]
    Euler's number as a \texttt{SimpleContinuedFraction} with integer
    part $2$ and body generator
    \[
      g(n) =
      \begin{cases}
        2\bigl(\lfloor(n+1)/3\rfloor + 1\bigr) & n \equiv 2 \pmod{3}, \\
        1 & \text{otherwise.}
      \end{cases}
    \]
    This encodes the classical pattern $e = [2;\,1,2,1,\,1,4,1,\,1,6,1,\ldots]$,
    where every third body coefficient (at positions $n = 2, 5, 8, \ldots$)
    equals $2, 4, 6, \ldots$ and the remaining coefficients are $1$.  The
    generator is stateless and computed in $O(1)$.

  \item[$\mathtt{phi}$]
    The golden ratio $\varphi = [1;\,\overline{1}]$ as a
    \texttt{PeriodicSimpleContinuedFraction} with integer part $1$ and
    period $(1)$.

  \item[$\mathtt{sqrt2}$, $\mathtt{sqrt3}$, $\mathtt{sqrt5}$]
    The PSCFs $[1;\,\overline{2}]$, $[1;\,\overline{1,2}]$,
    $[2;\,\overline{4}]$ for $\sqrt{2}$, $\sqrt{3}$, $\sqrt{5}$
    respectively.

  \item[$\mathtt{metallic\_mean}(n)$]
    Factory returning the $n$-th metallic mean
    $(\,n + \sqrt{n^2 + 4}\,)/2 = [n;\,\overline{n}]$ as a PSCF.  For $n = 1$
    this recovers $\varphi$; for $n = 2$ the silver mean, etc.
\end{description}

\begin{catenadef}[\texttt{catena.numbers.constants}]
\begin{lstlisting}[style=python]
from catena.numbers.constants import e, phi, sqrt2, sqrt3, sqrt5, metallic_mean

e.integer_part            # 2
e.generator(0)            # 1  (first body coefficient)
e.generator(2)            # 2  (a_3 = 2)
e.generator(5)            # 4  (a_6 = 4)

phi.period                # (1,)
sqrt2.period              # (2,)
sqrt3.period              # (1, 2)
metallic_mean(3).period   # (3,)   -- bronze mean = [3; 3, 3, ...]
\end{lstlisting}
\end{catenadef}

\subsection{Random SCFs and the Gauss--Kuzmin distribution}
\label{subsec:numbers-randoms}

\paragraph{Mathematical background (stated without proof).}
The theory of continued fractions has a natural probabilistic counterpart.
Equip $[0,1)$ with the Gauss measure $\mu$, whose density with respect to
Lebesgue measure is
\[
  \frac{d\mu}{d\lambda}(x) = \frac{1}{\ln 2} \cdot \frac{1}{1 + x}.
\]
Under $\mu$, the partial quotients $a_1, a_2, \ldots$ of a \emph{random}
continued fraction (drawn uniformly from $[0,1)$ with respect to
$\mu$) are independent and identically distributed with the
\emph{Gauss--Kuzmin distribution}:
\[
  \Pr[a_n = k] = \log_2\!\left(1 + \frac{1}{k(k+2)}\right),
  \qquad k \geq 1. \tag{GK}
\]
This result, due to Gauss (observed) and Kuzmin (proved), is the
foundational theorem of the ergodic theory of continued fractions.  A
full proof, and the connection to the Gauss--Kuzmin--Wirsing operator,
is deferred to a forthcoming section.

The Gauss--Kuzmin distribution can be sampled via the inverse-transform
method: if $U \sim \mathrm{Uniform}(0,1)$ then
\[
  K = \left\lfloor \frac{1}{2^U - 1} \right\rfloor \tag{IT}
\]
has the Gauss--Kuzmin distribution.  The library uses this formula
directly.

\paragraph{Implementation.}

\catena provides two concrete samplers and one abstract base class.

\begin{description}[leftmargin=2em, style=nextline]
  \item[\texttt{Seed}]
    A stateful SHA-256 chain: each access to \texttt{.state} returns the
    current bytes and advances the internal state by one hash round.
    The default global \texttt{\_SEED} is initialised from a random UUID,
    so different process invocations produce independent sequences.

  \item[\texttt{GaussKuzminSHA}]
    Stateless sampler for a fixed seed.  For index $n$, it derives a
    53-bit uniform integer via SHA-256 of $(\textit{seed} \| n)$, converts
    it to a float $U \in [0,1)$, and applies~(IT).  Each call is $O(1)$
    and the outputs for different $n$ are independent (under the
    cryptographic assumption that SHA-256 behaves as a random oracle).

  \item[\texttt{GaussKuzminSHAArbitrary}]
    As \texttt{GaussKuzminSHA} but the uniform sample $U$ is computed to
    an arbitrary number of significant decimal digits using
    \texttt{UniformSHAArbitrary} (a SHA-256 chain that assembles as many
    256-bit blocks as needed).  The result of~(IT) is then an exact Python
    \texttt{int}, avoiding any float truncation that could bias the floor
    computation for large partial quotients.

  \item[\texttt{RandomSCF} (abstract base class)]
    Defines the factory interface: \texttt{generator()},
    \texttt{cached\_generator()}, \texttt{finite\_generator(size)},
    \texttt{periodic\_generator(period\_size, pre\_period\_size)},
    \texttt{scf()}, \texttt{finite\_scf(size)}, \texttt{periodic\_scf(\ldots)}.
    Concrete subclasses supply \texttt{\_\_make\_callable\_\_()}, which
    returns a fresh callable backed by a new SHA seed drawn from the
    instance's \texttt{Seed}.

  \item[\texttt{GaussKuzminSCF}, \texttt{GaussKuzminArbitrarySCF}]
    Concrete subclasses of \texttt{RandomSCF} backed by
    \texttt{GaussKuzminSHA} and \texttt{GaussKuzminSHAArbitrary}
    respectively.
\end{description}

\begin{catenadef}[\texttt{catena.numbers.randoms}]
\begin{lstlisting}[style=python]
from catena.numbers.randoms import GaussKuzminSCF, GaussKuzminArbitrarySCF, Seed

# reproducible: fix the seed
rng = GaussKuzminSCF(seed=Seed(42))
scf = rng.scf(integer_part=0)        # infinite SCF with GK-distributed body
scf.convergent(10)                    # (p_10, q_10) -- some large rational

# finite sample of length 20
fscf = rng.finite_scf(size=20)       # FiniteSimpleContinuedFraction

# arbitrary-precision sampler
rng2 = GaussKuzminArbitrarySCF(precision=100, seed=Seed("hello"))
g = rng2.generator()
g(0)                                 # a positive integer drawn from GK

# periodic SCF with 5-term period, no pre-period
pscf = rng.periodic_scf(period_size=5, integer_part=3)
pscf.period                          # tuple of 5 positive integers
\end{lstlisting}
\end{catenadef}

\begin{remark}[Determinism and reproducibility]
  Because each call to \texttt{\_\_make\_callable\_\_()} draws a fresh
  SHA seed from the instance's \texttt{Seed}, two successive calls to
  \texttt{rng.generator()} produce \emph{different} (independent)
  generators.  A single \texttt{Seed} object therefore acts as a
  splittable PRNG: the sequence of seeds it emits is deterministic given
  the initial state, and each emitted seed drives a fully independent
  coefficient sequence.  Passing \texttt{seed=Seed(k)} with a fixed
  integer $k$ makes the entire computation reproducible.
\end{remark}

\begin{remark}[Forthcoming: measure theory and distribution of convergents]
  The Gauss--Kuzmin law~(GK) is the marginal distribution of any single
  partial quotient of a $\mu$-random continued fraction.  Several deeper
  results --- the exponential mixing of the Gauss map, the
  Gauss--Kuzmin--Wirsing spectral gap, Khintchine's theorem on the
  geometric mean of partial quotients, and the distribution of convergent
  denominators --- are not yet covered in this document and are planned
  for a forthcoming section on the ergodic theory of continued fractions.
\end{remark}
